\documentclass[11pt,a4paper]{article}
\usepackage{amsmath,amssymb,amsthm}
\usepackage{hyperref}
\usepackage{microtype}
\usepackage{textcomp}
\usepackage{doi}

\newtheorem{theorem}{Theorem}[section]
\newtheorem{lemma}[theorem]{Lemma}
\newtheorem{proposition}[theorem]{Proposition}
\newtheorem{corollary}[theorem]{Corollary}
\newtheorem{definition}[theorem]{Definition}
\theoremstyle{remark}
\newtheorem{remark}[theorem]{Remark}

\title{Quaternionic Rigidity of Admissible Morphisms:\\
Every Admissible $\Phi_{q,\rho}$ Necessarily Factors through $\mathfrak{su}(2)$}

\author{J\'er\^ome Beau}
\date{2026\\[2pt]\small Version 1.15.3}

\begin{document}
  \maketitle

  \begin{abstract}
    Let $\Phi_{q,\rho} \colon V_q \to V_\rho$ be a morphism from the Weil representation space
    of $\mathrm{Heis}_3(\mathbb{Z}/q\mathbb{Z})$ to a representation space of $2I \subset \mathrm{SU}(2)$,
    subject to three structural constraints imposed by the admissible projection $\Pi$:
    involution equivariance ($\chi \mapsto -\chi$), quadratic compatibility with the pair observable
    $\sigma_{\mathrm{pair}}$, and naturality with respect to admissibility.
    The main result of this paper is a \emph{rigidity theorem}: any morphism satisfying these three
    constraints is forced to factor through the three-dimensional real Lie algebra $\mathfrak{su}(2)$,
    identified with the admissible neutral traceless sector $\mathrm{Im}\,\mathbb{H} = \mathrm{Im}\,\Pi
    \cap \mathrm{Ntrl}$.
    The quaternionic structure is not a choice but the unique fixed point of the three constraints.
    This upgrades the result of O23~\cite{Beau2026a27}---which established $\dim_{\mathbb{R}}
    (\mathrm{Im}\,\mathbb{H}) = 3$---from a dimension count to a categorical universality statement:
    $\mathfrak{su}(2)$ is the unique minimal admissible target through which every $\Phi_{q,\rho}$
    must factor.
    This gives the pair-capacity exponent a representation-theoretic carrier in the minimal
    admissible non-abelian sector.  It does not derive the separate capacity-to-rate map:
    that prescription combines a changing-degree LPS equation with a fixed-degree Heisenberg
    observable and has no native carrier on the latter substrate~\cite{Beau2026sgn}.
    The canonical construction $\Phi_{q,\rho} = \rho \circ \iota \circ \pi$ is exhibited explicitly and
    proved unique up to unitary equivalence.
  \end{abstract}

  \paragraph{Keywords.}
    Emergence; non-injective projection; admissible morphisms; quaternionic rigidity; su(2); representation theory; Weil
    representation; pair observables; Born–Infeld constraint; spectral admissibility

    \clearpage
    \tableofcontents

    \clearpage
  \section{Introduction}

  \subsection{The problem: existence of an admissible morphism}

    Paper O26~\cite{Beau2026a30} identified the pair observable $\sigma_{\mathrm{pair}}(n) =
  \sigma_c(n) \cdot \sigma_{q-c}(n)$ as the Hilbert--Schmidt norm squared of a rank-one operator
    in $\mathrm{End}(V_\rho)$ for some representation $\rho$ of $2I \subset \mathrm{SU}(2)$.
    This identification rests on the existence of a morphism
    \[
      \Phi_{q,\rho} \colon V_q \longrightarrow V_\rho
    \]
    satisfying three structural constraints:
    \begin{enumerate}
      \item[(a)] \emph{Involution equivariance}: $\mathbb{Z}_2$-equivariance under $c \leftrightarrow
      q - c$, the discrete realisation of $\chi \mapsto -\chi$;
      \item[(b)] \emph{Quadratic compatibility}: $\sigma_{\mathrm{pair}}(n) \sim
      \|\Phi_{q,\rho}(v^{(n)}_c)\|^2 \cdot \|\Phi_{q,\rho}(v^{(n)}_{q-c})\|^2$;
      \item[(c)] \emph{Naturality}: independence from all pipeline choices; dependence on the
      admissible structure $(\Pi, S_{\mathrm{BI}}, Q_8)$ alone.
    \end{enumerate}
    Paper O26 stated the existence of such a $\Phi_{q,\rho}$ as Hypothesis~4.4.
    The present paper resolves this hypothesis.
    The resolution is stronger than mere existence: the three constraints force a unique
    factorisation, making $\Phi_{q,\rho}$ a \emph{structural necessity} rather than a canonical choice.

  \subsection{The answer: a rigidity theorem}

    The central result is the following.
    Any morphism $\Phi_{q,\rho}$ satisfying (a)--(c) must factor through the admissible quotient
    $H_{\mathrm{eff}} = \mathrm{Im}\,\Pi \cap \mathrm{Ntrl}$, and the quaternionic identification
    $H_{\mathrm{eff}} \simeq \mathfrak{su}(2)$ (O23~\cite{Beau2026a27}) is forced by the interplay of
    the three constraints.
    No other three-dimensional or higher-dimensional real structure is compatible with all three
    constraints simultaneously.
    The comparison is explicit (Table~\ref{tab:candidates}):

    \begin{table}[h]
      \centering
      \begin{tabular}{lcccl}
        \hline
        Structure                           & (a)        & (b)        & (c)        & Verdict            \\
        \hline
        Abelian $\mathbb{R}^n$              & \texttimes & \texttimes & \checkmark & excluded by O23    \\
        Standard $\mathbb{C}^n$             & \checkmark & \checkmark & \texttimes & not canonical      \\
        Arbitrary Hilbert space             & \checkmark & \checkmark & \texttimes & pipeline-dependent \\
        $\mathbb{H} \,/\, \mathfrak{su}(2)$ & \checkmark & \checkmark & \checkmark & unique candidate   \\
        \hline
      \end{tabular}
      \caption{Compatibility of candidate target structures with constraints~(a)--(c).}
      \label{tab:candidates}
    \end{table}

  \subsection{Relation to the representation-side O-series chain}

    O24~\cite{Beau2026a28} proves that vertical fibre structure does not perturb the
    observable rank.  The present result adds a representation-theoretic layer to the
    transfer-free implication chain
    \begin{align*}
      \text{emergence} &\Rightarrow \text{non-injectivity} \Rightarrow \text{pair structure}\\
      &\Rightarrow \text{quadratic form} \Rightarrow \mathfrak{su}(2),
    \end{align*}
    where each implication is a theorem of the O-series (the first being the no-go theorem of
    paper~L~\cite{Beau2026l}).
    The exponent $\delta_{\mathrm{pair}}$ thereby acquires a representation-theoretic carrier:
    it describes a Hilbert--Schmidt trajectory in the minimal admissible non-abelian sector
    $\mathfrak{su}(2)$.  No capacity-to-rate conclusion follows from this chain.

  \subsection{Scope}

    This paper proves the rigidity theorem (Section~\ref{sec:rigidity}) and the canonical
    construction (Section~\ref{sec:construction}).
    It does not execute the numerical tests of O26 (those require O25 data~\cite{Beau2026a29}).
    It does not posit or repair the LPS-to-Heisenberg capacity-to-rate transfer.

  \subsection{Organisation}

    Section~\ref{sec:setup} fixes notation and recalls the relevant results from O18--O24.
    Section~\ref{sec:exclusion} eliminates the abelian and generic Hilbert-space candidates.
    Section~\ref{sec:rigidity} states and proves the rigidity theorem.
    Section~\ref{sec:construction} exhibits the canonical factorisation and proves uniqueness.
    Section~\ref{sec:chain} states the representation-side result and the transfer boundary.
    Section~\ref{sec:open} states the remaining open problem.

  \section{Setup and Notation}
  \label{sec:setup}

  Let $q$ be an odd prime.
  The Weil representation of $\mathrm{Heis}_3(\mathbb{Z}/q\mathbb{Z})$ on $\mathbb{C}^q$
  decomposes over conjugate character pairs $\{c, q-c\}$ with $c \in (\mathbb{Z}/q\mathbb{Z})^\times$,
  yielding blocks $V_c \simeq \mathbb{C}^q$ and $V_{q-c} \simeq \overline{V_c}$.
  The pair space is $V_q := V_c \oplus V_{q-c}$.

  The parity involution $\chi \mapsto -\chi$, derived in O18~\cite{Beau2026a22} from the
  Born--Infeld evenness of the substrate action, acts on $V_q$ by complex conjugation:
  $V_{q-c} \simeq \overline{V_c}$.
  The canonical pair observable (O19~\cite{Beau2026a23}) is
  \[
    \sigma_{\mathrm{pair}}(n) := \sigma_c(n) \cdot \sigma_{q-c}(n),
    \qquad \sigma_c(n) = \delta r_n / |S_n|.
  \]

  The admissible neutral traceless sector is
  \[
    H_{\mathrm{eff}} := \mathrm{Im}\,\Pi \cap \mathrm{Ntrl},
  \]
  with $\dim_{\mathbb{R}} H_{\mathrm{eff}} = 3$ by O23 Theorem~3.2~\cite{Beau2026a27}, together
  with the quaternionic identification
  \[
    \iota \colon H_{\mathrm{eff}} \xrightarrow{\;\simeq\;} \mathfrak{su}(2),
  \]
  canonical and unique up to $\mathrm{SU}(2)$-conjugacy.
  The admissibility thread $Q_8 \subset 2I \subset \mathrm{SU}(2)$ selects a spin-$\tfrac{1}{2}$
  representation $\rho$ as the minimal non-abelian candidate (O26~\cite{Beau2026a30},
  Section~3.3).

  \begin{definition}[Naturality with respect to admissibility]
    \label{def:naturality}
    A morphism $\Phi \colon V_q \to W$ is \emph{natural with respect to admissibility} if it is
    invariant under all admissible endomorphisms of $V_q$, i.e.\ under every linear map
    $f \colon V_q \to V_q$ that preserves $\Pi$-admissibility in the sense that
    $f(\ker \Pi_{\mathrm{adm}}) \subseteq \ker \Pi_{\mathrm{adm}}$.
    Equivalently, $\Phi$ is natural if it depends only on the admissible structure
    $(\Pi, S_{\mathrm{BI}}, Q_8)$ and not on any auxiliary choice of basis, BFS block,
    or Gram--Schmidt frame.
  \end{definition}

  Throughout, $V_\rho$ denotes a complex representation space of $2I$ of dimension $d_\rho$,
  endowed with its natural Hilbert structure.

  \begin{remark}[Non-arbitrariness of $\rho$ via the admissibility thread]
    \label{rem:thread}
    The representation $\rho$ is not chosen freely.
    The admissibility thread
    \[
      Q_8 \;\subset\; 2I \;\subset\; \mathrm{SU}(2)
    \]
    is imposed structurally: $Q_8$ is the minimal non-abelian finite subgroup of $\mathrm{SU}(2)$
    compatible with bounded-flux admissibility~\cite{Beau2026a1}; $2I$ is the maximal admissible
    finite subgroup, distinguished by the isotropy of its second-moment tensor among all binary
    polyhedral groups~\cite{Beau2026a2}; and the LPS family provides the Ramanujan approximation
    to $\mathrm{SU}(2)$ in the large-$q$ limit~\cite{Beau2026a1}.
    The spin-$\tfrac{1}{2}$ sector is selected because $Q_8 \subset 2I$ acts on $V_{\rho_{1/2}}
    \simeq \mathbb{C}^2$ as the minimal non-abelian representation, yielding the smallest
    non-trivial Laplacian eigenvalue $\lambda_{\rho_4} = 6 < 8 = \lambda_{\mathrm{ab}}$ and hence
    the widest admissibility window~\cite{Beau2026a1,Beau2026a30}.
    The choice of $\rho$ is therefore an output of the admissibility hierarchy, not an input.
  \end{remark}

  \section{Elimination of Non-Quaternionic Candidates}
  \label{sec:exclusion}

  We show that no candidate target structure other than $\mathfrak{su}(2)$ is compatible with
  all three constraints.

  \begin{proposition}[Abelian structures are excluded]
    \label{prop:abelian}
    No abelian target $V_\rho$ satisfies conditions~(a), (b), and~(c) simultaneously.
  \end{proposition}

  \begin{proof}
    Condition~(c) forces $\Phi_{q,\rho}$ to factor through $H_{\mathrm{eff}} = \mathrm{Im}\,\Pi \cap
    \mathrm{Ntrl}$.
    By O23 Theorem~3.2~\cite{Beau2026a27}, $H_{\mathrm{eff}}$ is quaternionic and its real rank is
    exactly 3, maximal among all admissible associative structures compatible with the Born--Infeld
    parity.
    An abelian target would require $H_{\mathrm{eff}}$ to split as a direct sum of
    one-dimensional real factors, contradicting the non-commutativity of the quaternionic structure.
  \end{proof}

  \begin{proposition}[Generic complex Hilbert spaces are not canonical]
    \label{prop:hilbert}
    A morphism $\Phi_{q,\rho}$ targeting an arbitrary complex Hilbert space $\mathcal{H}$ can satisfy
    conditions~(a) and~(b) but cannot satisfy condition~(c) unless $\mathcal{H}$ admits a quaternionic
    structure compatible with admissibility.
  \end{proposition}

  \begin{proof}
    Conditions~(a) and~(b) can be satisfied by any space admitting a conjugation and a norm,
    hence in particular by any $\mathcal{H}$.
    Condition~(c) requires that $\Phi_{q,\rho}$ depend only on $(\Pi, S_{\mathrm{BI}}, Q_8)$.
    Any such morphism must factor through the admissible quotient $\pi \colon V_q \to H_{\mathrm{eff}}$
    (universality of the admissible quotient; see Lemma~\ref{lem:quotient}).
    The image $H_{\mathrm{eff}}$ is three-dimensional and quaternionic.
    A morphism $H_{\mathrm{eff}} \to \mathcal{H}$ that is natural with respect to admissibility must
    intertwine the $\mathrm{SU}(2)$-action, hence it is a morphism of representations.
    By Schur's lemma, such a morphism exists and is canonical only when $\mathcal{H}$ contains
    $V_\rho$ as an $\mathrm{SU}(2)$-subrepresentation, i.e.\ when $\mathcal{H}$ carries a
    compatible quaternionic structure.
  \end{proof}

  \begin{lemma}[Universality of the admissible quotient]
    \label{lem:quotient}
    Let $\pi \colon V_q \twoheadrightarrow H_{\mathrm{eff}} = V_q / \ker \Pi_{\mathrm{adm}}$ be the
    admissible quotient map.
    Any morphism $\Phi \colon V_q \to W$ satisfying condition~(c) factors uniquely through $\pi$:
    there exists $\bar{\Phi} \colon H_{\mathrm{eff}} \to W$ such that $\Phi = \bar{\Phi} \circ \pi$.
  \end{lemma}

  \begin{proof}
    Condition~(c) states that $\Phi$ depends only on the admissible structure.
    The kernel $\ker \Pi_{\mathrm{adm}}$ consists precisely of the directions invisible to
    admissible observables.
    Any morphism that is natural with respect to admissibility must vanish on $\ker
    \Pi_{\mathrm{adm}}$, for otherwise it would distinguish substrate configurations that are
    admissibility-equivalent.
    Vanishing on $\ker \Pi_{\mathrm{adm}}$ is exactly the factorisation condition through $\pi$.
  \end{proof}

  \section{Rigidity Theorem}
  \label{sec:rigidity}

  \begin{theorem}[Quaternionic rigidity]
    \label{thm:rigidity}
    Let $\Phi_{q,\rho} \colon V_q \to V_\rho$ be any morphism satisfying conditions~(a)--(c).
    Then:
    \begin{enumerate}
      \item[(i)] $\Phi_{q,\rho}$ factors through the admissible quotient $\pi \colon V_q \to
      H_{\mathrm{eff}}$;
      \item[(ii)] the intermediate space $H_{\mathrm{eff}}$ is necessarily isomorphic to $\mathfrak{su}(2)$
      as a real Lie algebra, via the quaternionic identification $\iota$ of
      O23~\cite{Beau2026a27};
      \item[(iii)] the target $V_\rho$ is necessarily an $\mathrm{SU}(2)$-representation space, and
      the morphism $H_{\mathrm{eff}} \to V_\rho$ is a representation morphism.
    \end{enumerate}
    In particular, under constraints~(a)--(c), any admissible morphism $\Phi_{q,\rho}$ necessarily
    factors through a three-dimensional real Lie algebra isomorphic to $\mathfrak{su}(2)$.
    The quaternionic structure is not a modelling choice but the unique fixed point of the
    three constraints.
    In particular, the result is independent of the specific Weil realisation of
    $\mathrm{Heis}_3(\mathbb{Z}/q\mathbb{Z})$ and holds for any admissible representation of the pair
    structure satisfying~(a)--(c), whatever the underlying group or prime~$q$.
  \end{theorem}

  \begin{proof}
    \textbf{Part~(i).}
    Follows immediately from Lemma~\ref{lem:quotient}: condition~(c) forces factorisation through
    $\pi$.

    \textbf{Part~(ii).}
    After factorisation, $\bar{\Phi} \colon H_{\mathrm{eff}} \to V_\rho$ must be natural with
    respect to the Born--Infeld admissibility.
    The only admissible automorphisms of $H_{\mathrm{eff}}$ are those compatible with the
    Born--Infeld parity (condition~(a)) and the quadratic norm (condition~(b)).
    By O23 Theorem~3.2~\cite{Beau2026a27}, the group of such automorphisms is $\mathrm{SU}(2)$,
    acting on $H_{\mathrm{eff}} \simeq \mathfrak{su}(2)$ by the adjoint action.
    No other Lie group structure on a three-dimensional real space is compatible with both
    the conjugation involution of~(a) and the positive-definite quadratic form of~(b) simultaneously;
    this is the classical uniqueness of $\mathbb{H}$ as the only normed associative real division
    algebra of dimension greater than~2 (Frobenius theorem).
    Hence $H_{\mathrm{eff}} \simeq \mathfrak{su}(2)$ is forced.

    \textbf{Part~(iii).}
    A natural morphism $\bar{\Phi} \colon \mathfrak{su}(2) \to V_\rho$ must intertwine the
    $\mathrm{SU}(2)$-action.
    By Schur's lemma, such a morphism exists and is non-zero only when $V_\rho$ contains an
    irreducible $\mathrm{SU}(2)$-subrepresentation; it is then a morphism of representations.
  \end{proof}

  \begin{remark}[Relation to O23]
O23~\cite{Beau2026a27} established that $\dim_{\mathbb{R}} H_{\mathrm{eff}} = 3$.
Theorem~\ref{thm:rigidity} upgrades this from a dimension statement to a \emph{categorical
universality} statement: $\mathfrak{su}(2)$ is the unique minimal admissible target through
which every morphism satisfying~(a)--(c) must factor.
The integer~3 is not an accident of the Weil representation: it is the unique dimension
compatible with the simultaneous constraints of involution, quadratic form, and
associativity required by naturality.
  \end{remark}

  \begin{remark}[Frobenius and admissibility]
The appeal to the Frobenius theorem in Part~(ii) is not circular.
Frobenius establishes that $\mathbb{R}$, $\mathbb{C}$, $\mathbb{H}$ are the only normed real
associative division algebras.
Dimension~1 ($\mathbb{R}$) is excluded by the non-commutativity imposed by~(c) and O23.
Dimension~2 ($\mathbb{C}$) is excluded because a complex structure on $H_{\mathrm{eff}}$ is
not preserved by the $\mathbb{Z}_2$-involution of~(a) (complex conjugation is not an inner
automorphism of $\mathbb{C}$).
Dimension~4 ($\mathbb{H}$ itself, not $\mathfrak{su}(2) \simeq \mathrm{Im}\,\mathbb{H}$) is
excluded because $\mathrm{dim}_{\mathbb{R}} H_{\mathrm{eff}} = 3$ (O23).
The unique compatible structure is therefore $\mathrm{Im}\,\mathbb{H} \simeq \mathfrak{su}(2)$.
The argument uses only the algebraic consequences of constraints~(a)--(c) and the
dimension bound of O23; it does not depend on the Heisenberg group structure or on
properties specific to any prime $q$.
  \end{remark}

  \section{Canonical Construction and Uniqueness}
  \label{sec:construction}

  \begin{definition}[Canonical admissible morphism]
    \label{def:Phi}
    Define
    \[
      \Phi_{q,\rho} \;:=\; \rho \circ \iota \circ \pi,
    \]
    where $\pi \colon V_q \to H_{\mathrm{eff}}$ is the admissible quotient map, $\iota \colon
    H_{\mathrm{eff}} \xrightarrow{\,\simeq\,} \mathfrak{su}(2)$ is the canonical quaternionic
    identification of O23, and $\rho \colon \mathfrak{su}(2) \to \mathrm{End}(V_\rho)$ is the
    Lie algebra action of the chosen irreducible $\mathrm{SU}(2)$-representation.
  \end{definition}

  \begin{theorem}[Existence, conditions, and uniqueness]
    \label{thm:main}
    The morphism $\Phi_{q,\rho}$ of Definition~\ref{def:Phi} satisfies conditions~(a)--(c).
    It is the unique morphism doing so, up to unitary equivalence: any other such morphism is of
    the form $U \circ \Phi_{q,\rho}$ for some unitary $U \in \mathrm{U}(V_\rho)$ commuting with
    the $Q_8$-action.
  \end{theorem}

  \begin{proof}
    \textbf{Conditions.}
    (a) The involution $\chi \mapsto -\chi$ acts on $\mathfrak{su}(2)$ by $X \mapsto -X$, hence on
    $V_\rho$ by $\rho(-X) = -\rho(X)$; the $\mathbb{Z}_2$-equivariance is preserved at each stage.
    (b) The norm $\|\Phi_{q,\rho}(v)\|^2 = \|\rho(\iota(\pi(v)))\|^2_{V_\rho}$ recovers the
    quadratic structure of $\sigma_{\mathrm{pair}}$ since $\pi$ preserves the norm class, $\iota$
    is an isometry up to a fixed scale, and $\rho$ is unitary.
    (c) Each factor $\pi$, $\iota$, $\rho$ depends only on $(\Pi, S_{\mathrm{BI}}, Q_8)$, by
    Lemma~\ref{lem:quotient}, O23, and the admissibility thread respectively.

    \textbf{Uniqueness.}
    By Theorem~\ref{thm:rigidity}, any $\Phi'_{q,\rho}$ satisfying (a)--(c) factors as
    $\rho' \circ \iota \circ \pi$ for some representation morphism $\rho'$.
    The ambiguity in $\iota$ is $\mathrm{SU}(2)$-conjugacy, which is unobservable in
    $\|\Phi_{q,\rho}(\cdot)\|^2$.
    The ambiguity in $\rho'$ versus $\rho$ is, by Schur's lemma, a unitary $U$ commuting with
    the $Q_8$-action.
  \end{proof}

  \section{Representation Result and Transfer Boundary}
  \label{sec:chain}

  Under Theorem~\ref{thm:main}, the pair observable satisfies
  \[
    \sigma_{\mathrm{pair}}(n) \;\sim\; \bigl\|\Phi_{q,\rho}(v^{(n)}_c) \otimes
    \Phi_{q,\rho}(v^{(n)}_{q-c})\bigr\|^2_{\mathrm{HS}},
  \]
  realising $\delta_{\mathrm{pair}}$ as the decay exponent of a Hilbert--Schmidt norm along a
  trajectory $(\widetilde{M}^{(n)})_{n \geq 0}$ in $\mathrm{End}(V_\rho)$.

  \begin{corollary}[Representation carrier of $\delta_{\mathrm{pair}}$]
    \label{cor:beta}
    The exponent $\delta_{\mathrm{pair}}$ is the decay exponent of a Hilbert--Schmidt
    trajectory in a representation space factoring through the minimal admissible
    non-abelian sector $\mathfrak{su}(2)$.
    This statement does not imply $\beta^*=1/(\delta_{\mathrm{pair}}+\tfrac12)$.
    That reciprocal requires a distinct cross-substrate growth hypothesis, refuted for the
    two native realisations in the frozen corpus~\cite{Beau2026sgn}.
  \end{corollary}

  \begin{remark}[Completeness of the chain]
The full implication chain is now:
\[
  \text{emergence} \Rightarrow \text{non-injectivity} \Rightarrow \text{pair structure}
  \Rightarrow \text{quadratic form} \Rightarrow \mathfrak{su}(2).
\]
The first three inputs are established respectively in paper~L, O18, and O26
\cite{Beau2026l,Beau2026a22,Beau2026a30}; the last arrow is
Theorem~\ref{thm:rigidity}.
Each arrow is a proved result of the O-series or companion paper programme.
No arrow is an assumption.
  \end{remark}

  \section{Remaining Open Problem}
  \label{sec:open}

  Theorem~\ref{thm:rigidity} establishes the structure of $\Phi_{q,\rho}$ unconditionally.
  The single remaining open question in the representation-theoretic component of the O-series is
  the numerical identification of the correct irreducible sector: which value of $d_\rho$ is
  realised by $\mathrm{O25}$ data?

  The spin-$\tfrac{1}{2}$ candidate ($d_\rho = 2$) predicts effective dimension $r_{\mathrm{eff}} =
  d_\rho^2 = 4$ in the covariance test of O26 Criterion~5.4~\cite{Beau2026a30}.
  A different fixed value $r_{\mathrm{eff}} = d_\rho^2 \neq 4$ would identify an alternative
  irreducible sector without invalidating the rigidity theorem, since
  Theorem~\ref{thm:rigidity} holds for any $\rho$.
  This test is executable on existing O25 data~\cite{Beau2026a29} without additional BFS
  computation.

  \section*{Acknowledgements}
  The author acknowledges the use of large language models as a supportive tool for refining
  language, structure, and internal consistency during the development of this manuscript.

  \bibliographystyle{plain}
  \bibliography{cosmochrony-bibliography}

\end{document}
